\documentclass[10pt,a4paper]{amsart}
\usepackage[margin=19mm]{geometry}
\usepackage{amsmath,amssymb,mathtools}
\usepackage[scr=rsfs,cal=euler]{mathalfa}
\usepackage{xfrac}
\usepackage[unicode,hidelinks,bookmarks=false]{hyperref}
\newcommand{\ie}{i.e.\ }
\newcommand{\cf}{cf.\ }
\newcommand{\resp}{respectively\ }
\newcommand{\Subsection}{\subsection}
\providecommand{\nobreakdash}{\nobreak\hbox{-}}
\setlength{\emergencystretch}{2em}
\multlinegap0pt
\usepackage{bm}
\usepackage{bbm}
\usepackage{dsfont}
\usepackage{xspace}
\usepackage{cancel}
\usepackage{mathtools}
\usepackage{amsthm}
\makeatletter
\@ifundefined{theorem}{\newtheorem{theorem}{Theorem}}{}
\@ifundefined{lemma}{\newtheorem{lemma}[theorem]{Lemma}}{}
\makeatother
\newcommand{\Vc}{{\mathcal V}_{X,G}}
\title[Totally mixed conditional independence equilibria of generic]{Totally mixed conditional independence equilibria of generic games}
\author{Matthieu Bouyer, Irem Portakal, Javier Sendra-Arranz}
\date{Source: arXiv:2511.11467v1 (CC BY 4.0)}
\begin{document}
\maketitle

\noindent\emph{Source and licence.} Totally mixed conditional independence equilibria of generic games. Version arXiv:2511.11467v1 (CC BY 4.0), distributed under the Creative Commons Attribution 4.0 International licence, as recorded in the arXiv licence metadata for this exact version and shown on the abstract page https://arxiv.org/abs/2511.11467v1. Full rights notice: licenses/arxiv-2511.11467-CC-BY-4.0.txt.\par\smallskip

\noindent\textit{Editorial scope.} Counted unit: the Theorem whose source label is \texttt{theo:emptyness} in the article (source0.tex lines 1534--1540, source label \texttt{theo:emptyness}), delivered together with the article's own complete original proof of it (source0.tex lines 1541--1608, one \texttt{proof} environment of 68 lines). No mathematics is written, repaired or re-derived: statement and proof are the article's own text line for line. Required context retained uncounted: eq:poly chow (lines 1294--1299); lemma: monomials (lines 1500--1502); lemma: inequality ds (lines 1519--1521). Every definition, display and label of those lines is retained unchanged. Omitted from this package and not redistributed: 1--1293, 1300--1499, 1503--1518, 1522--1533, 1609--2286. Nothing inside a retained excerpt is omitted or altered: every retained body is delivered exactly as the article's lines are written, so no omission receipt of any kind accompanies this package. Citations: the retained excerpts carry no citation command at all, so no bibliography entry of the article is transcribed and no reference is redistributed. Reference adaptation: none was needed and none was made; every reference inside the delivered unit resolves inside it, so no unresolved reference or citation is produced. Printed slips kept literal: no printed word, symbol, hypothesis or number of the counted statement or of its proof was altered. Layout and class: the article's own class is \texttt{amsart} with packages including geometry, xcolor, appendix, fontenc, amsthm, amsmath, hyperref, comment, caption, graphicx, subfigure, latexsym, amssymb, float; the wrapper is the lane's pinned \texttt{amsart} editorial wrapper at 10pt on A4, which restates the theorem-like environments the retained text needs and adds the article's own macro definitions that the retained text needs (\texttt{\textbackslash Vc}), with extra wrapper packages (bm, bbm, dsfont, xspace, cancel, mathtools). The delivered numbering is the unit's own. Assets: no retained span contains a figure environment, a graphics-inclusion command, a PDF-page inclusion, a table or a TikZ picture; no image file is copied into this package, the package declares no asset dependency and no image file is redistributed with it. Licence: version arXiv:2511.11467v1 of this article is distributed under the Creative Commons Attribution 4.0 International licence, as recorded in the arXiv licence metadata for this exact version and shown on the abstract page https://arxiv.org/abs/2511.11467v1. A full rights notice is delivered at licenses/arxiv-2511.11467-CC-BY-4.0.txt. Characters: every retained excerpt is pure ASCII, so the delivered engine, XeLaTeX with its default Latin Modern text font, reports no missing character.
    \begin{equation}
        \label{eq:poly chow}
        [\Vc]= \displaystyle\prod_{i\in\mathcal S}  \widehat{x_i}^{d_i-1} \,\prod_{i\not\in\mathcal S}\prod_{j=1}^{n_i}\left(    
        \sum_{l=0}^{d_{i,j}-1} x_i^l\left(\sum_{u=1}^kx_u\right)^{d_{i,j}-1-l}
        \right).
    \end{equation}

\begin{lemma}\label{lemma: monomials}
    The nonzero monomials that appear in the polynomial \eqref{eq:poly chow} are all the monomials $x_1^{\alpha_1}\cdots x_k^{\alpha_k}$ of degree $\displaystyle\sum_{l=1}^k S_l-n$ such that for $i\in\mathcal S$, $\displaystyle\alpha_i\leq \sum_{l=1}^kS_l -n+1-d_i$.
\end{lemma}

\begin{lemma}\label{lemma: inequality ds}
    For every $i\in[k]$, it holds that $S_i-n_i\le D_i-1$.
\end{lemma}

\begin{theorem}\label{theo:emptyness}
    Let $\mathcal S$ be the set of isolated vertices of $G$. Then, the Nash CI variety is nonempty for generic games if and only if
    \begin{equation}\label{eq: final nonempty}
    d_i\leq1+\frac12 \sum_{l=1}^k(D_l-1)
    \end{equation}
    for every $i\in\mathcal S$.
\end{theorem}

\begin{proof}
The Nash CI variety of a generic game $X$ is nonempty if and only if the class $\left[ \Vc\right]$ is nonzero. Thus, it is enough to show that the polynomial \eqref{eq:poly chow} is nonzero if and only if \eqref{eq: final nonempty} holds.
    To do so, we distinguish two cases: either  for all $ i\in\mathcal S$, $2d_i-2\leq \sum_{l=1}^k(S_l-n_l)$, or there exists $ i\in\mathcal S$ such that $2d_i-2>\sum_{l=1}^k(S_l-n_l)$.


Assume first that $2d_i-2\leq \sum_{l=1}^k(S_l-n_l)$ for all $i\in\mathcal S$. By Lemma \ref{lemma: inequality ds},  we deduce that \eqref{eq: final nonempty} holds. Therefore, in this case, to check the equivalence it is enough to check that the class \eqref{eq:poly chow} is nonzero.
Consider the monomial $\prod_{i=1}^kx_i^{S_i-n_i}$. This monomial has degree $\sum_{l=1}^k S_l-n$ and satisfies that for all $i\in\mathcal S$  $$\alpha_i=S_i-n_i=d_i-1\leq \sum_{l=1}^kS_l -n+1-d_i.$$ 
    Hence, by Lemma \ref{lemma: monomials}, this monomial appears in the polynomial \eqref{eq:poly chow}. Moreover,  such monomial is nonzero since $S_i-n_i\le D_i-1$ by Lemma \ref{lemma: inequality ds}. We conclude that the polynomial \eqref{eq:poly chow} is nonzero.


  Assume now that there exists $a\in\mathcal S$ such that $2d_a-2>\sum_{l=1}^k(S_l-n_l)$. Without loss of generality, we assume $d_a$ to be the maximum.  Then, for all $j\in\mathcal S\setminus \{a\}$, we have that 
    \begin{equation}\label{eq: unique da}
   2d_j-2\le d_j-1+d_a-1\le\sum_{l\in\mathcal S}(d_l-1)=\sum_{l\in\mathcal S}(S_l-n_l)\le\sum_{l=1}^k(S_l-n_l).
    \end{equation}
    Now, let $x^{\alpha}=x_1^{\alpha_1}\cdots x_k^{\alpha_k}$ be a monomial of the polynomial \eqref{eq:poly chow}. Then by Lemma \ref{lemma: monomials},
    \begin{equation}
        \label{eq:degree eq}
         \alpha_1+\cdots+\alpha_k=\sum_{l=1}^k S_l-n, \text{ and }\alpha_i\leq \sum_{l=1}^k S_l-n-d_i+1 \text{ for } i\in\mathcal S.
    \end{equation}
    Note that $x^\alpha$ is nonzero if and only if $\alpha_i\leq D_i-1$ for every $i$.
    Moreover, using \eqref{eq:degree eq} we get \begin{equation}\label{eq:alphas}
    \sum_{l\not\in\mathcal S}\alpha_l=\sum_{l=1}^k S_l-n-\sum_{l\in\mathcal S}\alpha_l=\sum_{l\not\in\mathcal S} (S_l-n_l)+\sum_{l\in\mathcal S}(d_l-1)-\sum_{l\in\mathcal S}\alpha_l.
    \end{equation}
    Assume first that \eqref{eq: final nonempty} holds. Then, for $i\in\mathcal S$, we fix the $i$--th coordinate of $\alpha$ to be
    \begin{equation}
        \label{eq: some alpha}
        \alpha_i=\min\big\{d_i-1,\sum_{l=1}^k S_l-n-d_i+1\big\}=\left\{\begin{array}{ll}
             d_i-1&\text{ if }i\ne a\text{ by \eqref{eq: unique da}}\\
             \displaystyle\sum_{l=1}^k S_l-n-d_i+1&\text{ if }i=a.
        \end{array}\right.
    \end{equation}
        Then, we get that  
        \begin{equation}\label{eq: sum of alphas}
\begin{array}{lll}
  \displaystyle  \sum_{l\not\in\mathcal S}\alpha_l&  \displaystyle =\sum_{l\not\in\mathcal S} (S_l-n_l)+\sum_{l\in\mathcal S}(d_l-1)-\sum_{l\in\mathcal S}\alpha_l&\text{ by \eqref{eq:alphas}}\\
  &  \displaystyle =\sum_{l\not\in\mathcal S} (S_l-n_l)+d_a-1-\big(\sum_{l=1}^k S_l-n-d_a+1\big)&\text{ by \eqref{eq: some alpha}}\\
  &  \displaystyle =2d_a-2-\sum_{l\in\mathcal S}(d_l-1).
\end{array}
        \end{equation}
        
\noindent Using \eqref{eq: final nonempty} and \eqref{eq: sum of alphas}, we deduce that
    \begin{equation*}
    \left\{\begin{array}{ll}
    \displaystyle\sum_{l\notin\mathcal S}\alpha_l\ge2d_a-2-\sum_{l=1}^k(S_l-n_l)>0\\
        \displaystyle\sum_{l\not\in\mathcal S}\alpha_l\le\sum_{l=1}^k(D_l-1)-\sum_{l\in\mathcal S}(d_l-1)=\sum_{l\notin\mathcal S}(D_l-1).
    \end{array}\right.
    \end{equation*}
    Hence, there exists a choice of $\alpha_i$ for $i\not\in\mathcal S$ such that $0\leq\alpha_i\leq D_i-1$. Since $\alpha_i\leq d_i-1=D_i-1$ for $i\in\mathcal S$ by \eqref{eq: some alpha}, we deduce that the polynomial \eqref{eq:poly chow} is nonzero since its monomial $x^\alpha$ is a non zero monomial.

Conversely, assume that \eqref{eq: final nonempty} does not hold. Then, \begin{equation}\label{eq:equivalent to final criterion}
        \displaystyle\sum_{l=1}^k(D_l-1)<2d_a-2.
    \end{equation}
    We claim that the polynomial \eqref{eq:poly chow} vanishes. Let $x^\alpha$ be a monomial of \eqref{eq:poly chow}.
    Note that the monomial $x^\alpha$ is zero if there exists $i\in\mathcal S$ such that $\alpha_i\ge d_i$. Hence, we may assume that $\alpha_i\leq d_i-1$ for every $i\in\mathcal S$. By Lemma \ref{lemma: monomials}, we may assume that 
    \begin{equation}\label{eq:ineq alphas}
    \alpha_i\leq \min\{d_i-1,\sum_{l=1}^k S_l-n-d_i+1\}=\left\{\begin{array}{ll}
             d_i-1&\text{ if }i\ne a\text{ by \eqref{eq: unique da}}\\
             \displaystyle\sum_{l=1}^k S_l-n-d_i+1&\text{ if }i=a
        \end{array}\right.
    \end{equation}
    for all $i\in\mathcal S$. Using this inequality we deduce that 
    \begin{align*}
        \sum_{l\not\in\mathcal S}\alpha_l&=\sum_{l\not\in\mathcal S} (S_l-n_l)+\sum_{l\in\mathcal S}(d_l-1)-\sum_{l\in\mathcal S}\alpha_l&\text{ by \eqref{eq:alphas}}\\
        &\geq \sum_{l\not\in\mathcal S} (S_l-n_l)+d_a-1-\big(\sum_{l=1}^kS_l-n-d_a+1\big)&\text{ by \eqref{eq:ineq alphas}}\\
        &=2d_a-2-\sum_{l\in\mathcal S}(d_l-1)>\sum_{l\notin\mathcal S}(D_l-1).&\text{ by \eqref{eq:equivalent to final criterion}}
    \end{align*}
    By Pigeonhole Principle, there exists $i\not\in\mathcal S$ such that $\alpha_i>D_i-1$, and the monomial $x^\alpha$ vanishes. We conclude that the polynomial \eqref{eq:poly chow} is zero.
\end{proof}

\end{document}
